\documentclass[oneside]{memoir}
\usepackage{ctex}
\usepackage{xcolor,calc}
\usepackage[T1]{fontenc}
\usepackage{kpfonts}
\usepackage{minted2}
\usepackage{tikz}
\usepackage{xstring}
\usepackage{indentfirst}
\usepackage[colorlinks=true, linkcolor=blue]{hyperref}
\usepackage{cleveref}

\usetikzlibrary{positioning}

% Cleveref configuration
\crefformat{chapter}{第~#2#1#3~章}
\crefformat{section}{~#2#1#3~节}
\crefformat{subsection}{~#2#1#3~节}

% Memoir chapter format, taken from:
% https://mirror.ox.ac.uk/sites/ctan.org/info/latex-samples/MemoirChapStyles/MemoirChapStyles.pdf
\setSingleSpace{1.1}
\SingleSpacing
\definecolor{chaptercolor}{gray}{0.8}
% helper macros
\newcommand\numlifter[1]{\raisebox{-2cm}[0pt][0pt]{\smash{#1}}}
\newcommand\numindent{\kern37pt}
\newlength\chaptertitleboxheight
\makechapterstyle{hansen}{
\renewcommand\printchaptername{\raggedleft}
\renewcommand\printchapternum{%
\begingroup%
\leavevmode%
\chapnumfont%
\strut%
\numlifter{\thechapter}%
\numindent%
\endgroup%
}
\renewcommand*{\printchapternonum}{%
\vphantom{\begingroup%
\leavevmode%
\chapnumfont%
\numlifter{\vphantom{9}}%
\numindent%
\endgroup}
\afterchapternum}
\setlength\midchapskip{0pt}
\setlength\beforechapskip{0.5\baselineskip}
\setlength{\afterchapskip}{3\baselineskip}
\renewcommand\chapnumfont{%
\fontsize{4cm}{0cm}%
\bfseries%
\sffamily%
\color{chaptercolor}%
}
\renewcommand\chaptitlefont{%
\normalfont%
\huge%
\bfseries%
\raggedleft%
}%
\settototalheight\chaptertitleboxheight{%
\parbox{\textwidth}{\chaptitlefont \strut bg\\bg\strut}}
\renewcommand\printchaptertitle[1]{%
\parbox[t][\chaptertitleboxheight][t]{\textwidth}{%
%\microtypesetup{protrusion=false}% add this if you use microtype
\chaptitlefont\strut ##1\strut}%
}
}
\chapterstyle{hansen}
\aliaspagestyle{chapter}{empty}

% Adjust memoir margins
\setlrmarginsandblock{3cm}{3cm}{*} % left/right margins
\setulmarginsandblock{2cm}{2cm}{*} % top/bottom margins
\checkandfixthelayout
% Adjust memoir title
\pretitle{\begin{center}\Huge\bfseries}
\posttitle{\par\vskip 2cm\end{center}}
\preauthor{\begin{center}\Large}
\postauthor{\par\end{center}}
\predate{\begin{center}\large}
\postdate{\par\end{center}}
\renewcommand{\ttdefault}{lmtt}

\begin{document}

\title{操作系统设计文档}
\author{黄越}
\date{\today}
\maketitle
\newpage

\tableofcontents

\chapter{前言}
\label{chap:preface}

这是我的第一个 OS 内核。既然没有经验，看不懂别人的操作系统，只好自己从零搭一个了。这个操作系统没有使用任何来自第三方的代码，不过使用了 OS 中常见的概念和抽象，例如 VFS 等。或许很多东西并不是最佳实践，但是能跑起来才是最重要的。

从项目数据上来看，\input{data/git_summary.tex}。请注意这里的“代码行数”包括所有被 git 统计的文件，例如这篇文档也在其中：LaTeX 代码也是代码。

从语言上来看，这个 OS 完全采用 C++ 实现。请注意这是 C++20，并且采用了许多 gcc/clang 独有的扩展；微软的 MSVC 或旧版本的编译器是无法编译它的。这是一个与几乎所有现有操作系统都不相同的决定。在\cref{chap:why_cpp}中，我介绍了我利用的许多 C++ 特性，以及在没有 C++ runtime 时的一些 workaround。希望在阅读这一章后，你可以理解我选择它的原因。

在这篇文档中，我将以 RISC-V 作为例子，介绍操作系统的启动流程以及各个子系统的实现方式。LoongArch 的支持准备由条件编译指令 \mintinline{cpp}{#ifdef __loongarch__} 控制，但我尚未开始。

这篇文档在一定程度上也算是我的 OS 学习笔记，因此语言上不会非常正式。此外，我并不清楚一部分术语的中文翻译，所以将会用对应的英文词汇替代它们。

\chapter{启动流程}
\label{chap:boot}
% 启动流程

\section{内存布局}

这个 OS 的 \textit{Load Memory Address} (LMA) 如下图所示。这是它们被加载到硬件上时的物理内存位置，但并非虚拟内存位置。换句话说，虚拟内存位置 (VMA) 是给链接器看的，而 LMA 则是给加载器看的。

\newcount\foreachcount
\foreachcount=1
\begin{center}
\begin{tikzpicture}[node distance=0, every node/.style={minimum width=3.5cm, minimum height=1cm, align=center}]
  \foreach \name/\label/\ptr in {
    textlow/.text.low/0x8020'0000,
    ptroot/pt\_root/0x8020'1000,
    a0/a0 (hart id)/0x8020'2008,
    a1/a1 (FDT address)/0x8020'2010,
    text2/.text/0x8020'3000,
    rodata/.rodata/,
    data/.data/,
    bss/.bss/,
    stack/kernel stack/stack bottom
} {
    \ifnum\foreachcount=1
      \node (\name) [draw] {\label};
    \else
      \node (\name) [draw, below=of \prev] {\label};
    \fi

    % Text-only pointer node
    \IfEq{\ptr}{}{}{
      % if (ptr != "")
      \node (p\name) [minimum width=1cm, right=4cm of \name, anchor=south] {\ptr};
      \draw[->, thick] (p\name.west) -- ([xshift=1mm]\name.north east);
    }

    % Keep track of previous node
    \xdef\prev{\name}
    \global\advance\foreachcount by 1
  }

  % The stack top
  \node (last) [minimum width=1cm, right=4cm of stack, anchor=north] {. + 131072};
  \draw[->, thick] (last.west) -- ([xshift=1mm]stack.south east);
\end{tikzpicture}
\end{center}

其中，pt\_root 是页表的根。这个内存布局要求可以由链接器脚本 link.ld 完成。

从虚拟内存的角度而言，.text.low、pt\_root 以及 a0 与 a1 的 VMA 就等于它们的 LMA。其余部分的 VMA 都是 LMA 增加一个特定的偏移量 0xffff'ffc0'0000'0000 所得到的。这主要是因为 RISC-V 指令集\cite{rvPrivManual}的要求。

我采用的 Sv39 页表有 39 位的虚拟地址，但是放入的是 64 位寄存器。RISC-V 要求前 25 位必须 sign-extend，因此虚拟地址空间被强制分为了两部分。这个 OS 将会占据较高的半部分虚拟地址，而 user-space 程序将会使用较低的一半。上面的偏移量的前 26 位是 1，正是这个原因。

实际上，链接器——至少 GNU 的 ld——会将 VMA 按照 8-byte 对齐，但却不会对 LMA 这么做。这会导致一些让人摸不着头脑的 bug，因此我们需要在链接器脚本里通过 VMA 计算 LMA，而不是反过来。

\section{启动}

在启动时，我决定立即启用虚拟内存，而不先进行其他任何设置。甚至寄存器 sp 都没有被设置，因此栈也暂时无法使用。启动时的 a0 和 a1 是 hart id 和 FDT (\textit{flattened device tree}) 的地址，也需要顺手保存到上面图中的位置。

我先按照偏移量，开 16 个 1 GB 的页，将前 16 GB 物理内存映射至对应的虚拟地址。（QEMU 本身只有 128 MB 的物理内存，虽然理论而言这读了设备树才能知道。）接下来，我将第 8 GB 恒等映射，然后启用页表。这个恒等映射是必要的，不然启用页表后的下一条指令将会立即 page-fault。之后我会把它移除。

接下来，OS 将会跳转到真正的 .text 段开始执行。普通的 jal 是不可以的，因为指令集编码限制，它只支持 20 位的跳转，而我的跨度有 64 位那么大。同样地，la 也是不可以的——它实际上是 auipc，也只支持 20 位范围内的寻址。因此，我只能将高段的 .text 固定下来，让它留在 0x8020'3000 的位置。

接下来，我会启动一个 free-list allocator。它是一个物理分配器，而且只管理 16 MB 内存。这是为了让大部分 C++ 容器可以使用，请见 \cref{sec:container}。

下面我将启动 PLIC，接受硬件中断。不过这需要分为几步。



\chapter{子系统}
\label{chap:subsystem}
% interrupt handler
\section{中断处理}
\label{sec:interrupt}

每个进程都有两个栈：一个是给自己的，一个是给 kernel 的。

在中断的时候，我们想把所有的寄存器都存到 kernel 的栈上，不过我们显然需要一个额外的临时寄存器。RISC-V 提供了一个名为 sscratch 的 CSR，我选择用它保存 user-stack 的位置。这就是 PCB 中的 usp 字段（\cref{sec:pcb}）。

这样我们就腾出了一个 sp 寄存器。我们需要在只有这一个寄存器的情况下找到当前进程 kernel-stack 的位置，也就是 PCB 中的 ksp 字段。在 C++ 中，这是 \cpp{scheduler.active->ksp}，所以只需要访问两次成员变量就可以了。我们使用 \cpp{static_assert} 来保证偏移量和汇编中的偏移量能对上。此外，VSCode 的 clangd 插件是可以直接显示 class/struct 中每个字段的偏移量和 padding 的，因此获取这个值也不是难事。

接下来，我们把 31 个寄存器存进 ksp 中，顺便存好 sscratch, sepc 和 sstatus。（显然寄存器 zero 是无需保存的。）我们还要靠它们返回中断发生的位置。

下面我们就可以调用 C++ 的中断处理器了。它可能会更改 \cpp{scheduler.active}，所以从中断中返回时还需要重新获取当前进程 ksp 的位置。

最后，重新载入 31 个寄存器以及 sepc 和 sstatus，然后把 sp 设置为 sscratch，就完成处理了。

\subsection{中断处理器}
\label{sec:interrupthandler}

这其实就是一个大型 switch。它可以依靠 ksp 读取寄存器的值，而 syscall 的参数就依赖这个实现。

除此以外似乎没什么值得说的；各种 syscall 的处理方式已经在相关的章节中描述过了。


\chapter{C++}
\label{chap:why_cpp}
% 介绍 C++
我大量使用了 C++ 的语言特性。即使没有 C++ runtime，这个语言依旧强大。接下来，我将描述我所使用的各种特性，以及对于缺失 runtime 时的一些解决方案。

在我看来，操作系统使用 C 主要是历史原因，而 C++ 不失为一种好的选择。至于 Rust，我个人感觉它“不够自由”，所以我不喜欢它。

\section{与汇编的交流}

众所周知，只要使用 \mintinline{cpp}{extern "C"}, 就可以将 C++ 函数当做 C 链接。但这不总是可行：大部分函数都被放在了 \mintinline{cpp}{namespace os} 之中，所以必须进行 name-mangling。好在它的规则并不复杂，只是初看吓人而已。具体的规则可以参见 Itanium C++ ABI \cite{itanium}, Section 5.1。如果熟悉 clang （前端）的内部实现，那么对它应该不会陌生。

举个例子，如\cref{sec:interrupt}所述，在 interruption handler 中，我们会加载 ksp。在汇编中，我们有如下代码：
\begin{minted}{asm}
stvec_pos:
  # Record current sp.
  csrw sscratch, sp
  
  # Load the ksp of the process.
  la sp, _ZN2os9schedulerE # os::scheduler
  ld sp, 48(sp)            # scheduler_t::active
  ld sp, 32(sp)            # pcb_t::ksp

  # Save registers to the kernel stack.
  sd ra, 0(sp)

  ...
\end{minted}

其中的 \_ZN2os9schedulerE 就是全局变量 os::scheduler 的符号名称，而后面的两条指令则是在加载它的成员变量。我们可以使用 static\_assert 来保证偏移量正确。

\section{容器}
\label{sec:container}


\bibliographystyle{plain}
\bibliography{bibtex}

\end{document}
